
\begin{equation}\tag{3.3.1}
 f^*:\Grb(X)\longrightarrow\Grb(X').
\end{equation}

\subsection*{3.4.}
Supposons enfin que les morphismes (3.2.3) et (3.3.1) soient tous les deux définis pour un morphisme \(f:X'\to X\) tel que \(r=d\). C'est le cas par exemple si \(X\) est quasi-compact avec un \(\underline{O}_X\)-module ample et si \(X'=P_X(\underline{E})\), où \(\underline{E}\) est un \(\underline{O}_X\)-module localement libre de rang \(r+1\). Alors la formule de projection dans la catégorie dérivée donne, pour tout \(x\in\Filt^i(X')_{\mathbb Q}\) et \(y\in\Filt_j(X)_{\mathbb Q}\), la formule de projection
\[
 f_*\bigl(x\cdot f^*(y)\bigr)=f_*(x)\cdot y\in\Filt_{j+r}(X)_{\mathbb Q}.
\]
Lorsque \(X\) est noethérien, alors 1.3.2 permet d'en déduire la formule graduée correspondante.

\section*{4. Nombres d'intersection}

\subsection*{4.1.}
Considérons un schéma propre \(X\) sur un corps \(k\). En composant le produit (3.1.1)
\[
 \Gr^i(X)\times\Gr_i(X)\longrightarrow \Gr_0(X)=\Filt_0(X)
 \qquad (i\in\mathbb Z)
\]
avec la caractéristique d'Euler--Poincaré
\[
 \chi:\Filt_0(X)\longrightarrow\mathbb Z,
\]
on obtient un accouplement appelé «~\emph{nombre d'intersection}~» :
\begin{equation}\tag{4.1.1}
 \Gr^i(X)\times\Gr_i(X)\longrightarrow \mathbb Z,
\end{equation}
